\chapter{Numerical series, products, and rearrangements}
\SourceRef{Rudin, Chapter 3 (including Theorems 3.33--3.55); Apostol, Chapter 8; Tao, Analysis I, Chapter 7. The proofs below use completeness and integer roots, not elementary transcendental functions. The explicitly identified logarithmic applications occur only after Chapter 7's constructions.}

\KEntry{K03.11.partial_sums}{Series and finite changes}
\begin{definition}
For scalars $a_n\in\R$ or $\C$, indexed by $n\geq0$, put $s_N=\sum_{n=0}^N a_n$. The equation $\sum_{n=0}^\infty a_n=S$ means $s_N\to S$ with $S$ finite. A divergent series has no such scalar sum.
\end{definition}
\begin{proof}
If $b_n=a_n$ for $n>m$, then for $N\geq m$, $\sum_{n=0}^N b_n=s_N+\sum_{n=0}^m(b_n-a_n)$. Thus either both series converge or neither does; their sums differ by the displayed finite correction. The limit laws similarly prove linearity of convergent series.
\end{proof}

\KEntry{K03.11.necessary}{Term test and Cauchy criterion}
\begin{theorem}
A series converges if and only if for every $\eps>0$ there is $N$ such that $\abs{\sum_{n=p}^q a_n}<\eps$ whenever $q\geq p\geq N$. Convergence implies $a_n\to0$, whose converse is false.
\end{theorem}
\begin{proof}
The finite tail is $s_q-s_{p-1}$, so the assertion is precisely the Cauchy criterion in the complete scalar field (\KRef{K03.01.cauchy}). Also $a_n=s_n-s_{n-1}\to0$. For the converse counterexample, $a_n=1/n$ tends to zero, whereas the terms from $2^k$ through $2^{k+1}-1$ sum to at least $1/2$; its partial sums are unbounded.
\end{proof}

\KEntry{K03.11.nonnegative}{Nonnegative sums and regrouping}
\begin{theorem}
For $a_n\geq0$, the series converges exactly when its partial sums are bounded. Otherwise write $\sum a_n=+\infty$. Its extended sum is $\sup\{\sum_{n\in F}a_n:F\subset\N_0\text{ finite}\}$, and is unchanged by any ordering or by a partition of the indices into finite or countable blocks, each summed nonnegatively.
\end{theorem}
\begin{proof}
The partial sums increase, so completeness gives their limit, or unboundedness gives divergence to $+\infty$. Every finite set lies in an initial segment and each initial segment is a finite set, proving the supremum formula. Every finite selection within finitely many blocks is a finite selection of original terms, giving one inequality for iterated sums. Conversely a finite set of original indices meets only finitely many blocks, whose sums dominate its sum. Taking suprema proves the reverse inequality, also when the supremum is infinite.
\end{proof}

\KEntry{K03.11.comparison}{Comparison tests}
\begin{theorem}
For eventually nonnegative $a_n,b_n$, $a_n\leq Cb_n$ eventually implies convergence of $\sum a_n$ if $\sum b_n$ converges. The equivalent divergence implication follows by contraposition. If $b_n>0$ eventually and $a_n/b_n\to L\in(0,\infty)$, the series converge together. If $L=0$, convergence of $\sum b_n$ implies that of $\sum a_n$; if $L=+\infty$, convergence of $\sum a_n$ implies that of $\sum b_n$.
\end{theorem}
\begin{proof}
Remove a finite prefix and compare partial sums. For $0<L<\infty$ use $L/2\leq a_n/b_n\leq2L$ eventually. For $L=0$ use $a_n\leq b_n$ eventually; for $L=\infty$ use $b_n\leq a_n$. No reverse implication in these last cases follows: $a_n=1/n^2,b_n=1/n$ gives a zero ratio and unequal behavior; $a_n=1/n,b_n=1/n^2$ gives an infinite ratio and unequal behavior.
\end{proof}

\KEntry{K03.11.geometric_p}{Geometric, harmonic, and p-series}
\begin{example}
The geometric series $\sum_{n\geq0}q^n$ has sum $(1-q)^{-1}$ for $\abs q<1$ and diverges for $\abs q\geq1$. The harmonic series diverges. For real $p$, $\sum_{n\geq1}n^{-p}$ converges exactly when $p>1$.
\end{example}
\begin{proof}
The finite identity $(1-q)\sum_{n=0}^Nq^n=1-q^{N+1}$ proves the geometric assertion: $\abs q^n\to0$ for $\abs q<1$ by monotone convergence (a positive limit would equal $\abs q$ times itself); otherwise the term test applies. Harmonic divergence was proved in \KRef{K03.11.necessary}. For $p\leq0$, terms do not tend to zero. For $0<p\leq1$, compare with $1/n$. For $p>1$, condensation, proved below without using this example, gives the geometric series $\sum_k2^{k(1-p)}$, with ratio $2^{1-p}<1$. Here the meaning and order laws of arbitrary real powers are the downstream interface \KRef{K03.07.real_powers}; for positive integral $p$ this proof needs only arithmetic. Nothing in the root or general power-series test uses this real-power application.
\end{proof}

\KEntry{K03.11.ratio}{Ratio test}
\begin{theorem}
For eventually nonzero terms, if $\limsup\abs{a_{n+1}/a_n}<1$, then $\sum a_n$ converges absolutely. If $\abs{a_{n+1}/a_n}\geq1$ eventually, the series diverges; in particular $\liminf\abs{a_{n+1}/a_n}>1$ suffices. A ratio limit $1$ is inconclusive.
\end{theorem}
\begin{proof}
Choose $q<1$ strictly above the limsup; for $n\geq N$, $\abs{a_n}\leq\abs{a_N}q^{n-N}$, whose sum is geometric. Absolute convergence implies convergence by \KRef{K03.11.absolute}, whose proof is independent of this test. In the second case the eventually nonzero magnitudes never decrease, so cannot tend to zero. The series $\sum1/n$ and $\sum1/n^2$ have ratio limit $1$ and different behavior (for the latter use condensation below).
\end{proof}

\KEntry{K03.11.root}{Root test without logarithms}
\begin{theorem}
Let $L=\limsup_{n\to\infty}\abs{a_n}^{1/n}\in[0,\infty]$, using positive integer roots from \KRef{K03.01.positive_roots}. If $L<1$, the series converges absolutely. If $L>1$, it diverges. When $L=1$, no conclusion follows.
\end{theorem}
\begin{proof}
For $L<q<1$, eventually $\abs{a_n}^{1/n}\leq q$, hence $\abs{a_n}\leq q^n$. The finite identity $\sum_{n=p}^m t^n=t^p(1-t^{m-p+1})/(1-t)$ for $0\leq t<1$, with $t^p\to0$ by monotone convergence, makes the comparison tails arbitrarily small. Absolute convergence follows directly from the numerical Cauchy criterion. For $1<q<L$, infinitely many $n$ satisfy $\abs{a_n}>q^n\geq1$, so the term test fails. For inconclusiveness compare $1/n$ and $1/n^2$: their roots tend to $1$. Indeed, for any $h>0$, the binomial theorem gives $(1+h)^n\geq\binom n3 h^3>n^2$ eventually, while $1\leq n^{2/n}$; squeeze and reciprocation prove the claim. No exponential or logarithm has been used.
\end{proof}

\KEntry{K03.11.condensation}{Cauchy condensation}
\begin{theorem}
If $a_n\geq0$ is nonincreasing for $n\geq1$, then $\sum a_n$ converges if and only if $\sum_{k\geq0}2^k a_{2^k}$ converges.
\end{theorem}
\begin{proof}
On a dyadic block,
\[2^k a_{2^{k+1}}\leq\sum_{n=2^k}^{2^{k+1}-1}a_n\leq2^k a_{2^k}.\]
Summing these inequalities compares block partial sums with the condensed partial sums, the left bound being one half of the condensed sum with its first term removed. Boundedness of one sequence of partial sums is therefore equivalent to boundedness of the other. Nonnegativity controls all intermediate partial sums.
\end{proof}

\KEntry{K03.11.alternating}{Alternating series and remainders}
\begin{theorem}
If $b_n\geq0$, $b_{n+1}\leq b_n$, and $b_n\to0$, then $\sum_{n\geq0}(-1)^nb_n$ converges. The remainder after $n=N$ has sign $(-1)^{N+1}$ or is zero, and magnitude at most $b_{N+1}$.
\end{theorem}
\begin{proof}
The partial sums ending at odd indices increase because $b_{2j}-b_{2j+1}\geq0$, while those ending at even indices decrease because $-b_{2j+1}+b_{2j+2}\leq0$. Each odd partial sum is below each later even partial sum, and the gap between adjacent such sums is a term $b_n\to0$. The two bounded monotone sequences thus have the same limit. Applying the same bracketing to the tail places its signed value between $0$ and $b_{N+1}$.
\end{proof}

\KEntry{K03.11.summation_parts}{Finite summation by parts}
\begin{theorem}
For $A_n=\sum_{j=0}^n a_j$, $A_{-1}=0$, and $0\leq p\leq q$,
\[\sum_{n=p}^q a_nb_n=A_qb_q-A_{p-1}b_p+\sum_{n=p}^{q-1}A_n(b_n-b_{n+1}).\]
Equivalently, with $T_n=\sum_{j=p}^n a_j$, the right side is $T_qb_q+\sum_{n=p}^{q-1}T_n(b_n-b_{n+1})$.
\end{theorem}
\begin{proof}
Insert $a_n=A_n-A_{n-1}$ and shift the index in the second finite sum. The uncancelled endpoints are exactly $A_qb_q$ and $-A_{p-1}b_p$. Replacing $A_n$ by $T_n$ gives the second formula. Limits may be taken only after the resulting boundary and series terms have been controlled.
\end{proof}

\KEntry{K03.11.dirichlet}{Dirichlet's test}
\begin{theorem}
If partial sums of scalar $a_n$ are bounded and real $b_n$ is monotone with limit zero, then $\sum a_nb_n$ converges.
\end{theorem}
\begin{proof}
Suppose first $b_n\downarrow0$ and $\abs{A_n}\leq M$. Then $\abs{T_n}\leq2M$, and summation by parts bounds every tail by
\[2M\left(b_q+\sum_{n=p}^{q-1}(b_n-b_{n+1})\right)=2Mb_p\longrightarrow0.\]
Use the Cauchy criterion. An increasing sequence tending to zero is nonpositive; replace it by its negative.
\end{proof}

\KEntry{K03.11.abel}{Abel's test}
\begin{theorem}
If $\sum a_n$ converges and real $b_n$ is monotone and bounded, then $\sum a_nb_n$ converges.
\end{theorem}
\begin{proof}
Monotone convergence gives $b_n\to b$. Split $a_nb_n=ba_n+a_n(b_n-b)$. The second factor in the latter product is monotone and tends to zero; the partial sums of $a_n$ are bounded, so Dirichlet applies.
\end{proof}

\KEntry{K03.11.absolute}{Absolute and conditional convergence}
\begin{theorem}
If $\sum\abs{a_n}<\infty$, then $\sum a_n$ converges, and every convergent tail satisfies $\abs{\sum_{n=p}^\infty a_n}\leq\sum_{n=p}^\infty\abs{a_n}$. Convergence without absolute convergence is called conditional.
\end{theorem}
\begin{proof}
For finite tails the inequality is the triangle inequality. Its right side tends uniformly to zero as $p\to\infty$, so the scalar Cauchy criterion gives convergence. Passing $q\to\infty$ in the finite inequality gives the displayed bound. The alternating harmonic series is conditional by the alternating test and harmonic divergence.
\end{proof}

\KEntry{K03.11.rearrange_absolute}{Grouping and absolute rearrangements}
\begin{theorem}
Consecutive finite grouping of a convergent series preserves its sum. Every bijective rearrangement of an absolutely convergent scalar series is absolutely convergent with the same sum.
\end{theorem}
\begin{proof}
Grouped partial sums are a subsequence of the original partial sums. For a permutation $\pi$, choose $m$ so $\sum_{n>m}\abs{a_n}<\eps$. Once a rearranged partial sum contains all indices $0,\ldots,m$, its difference from the original sum has magnitude at most $\sum_{n>m}\abs{a_n}<\eps$: subtract the finite included terms and take limits in the remaining absolutely convergent tail. The nonnegative regrouping theorem applied to $\abs{a_n}$ proves absolute convergence of the rearrangement.
\end{proof}

\KEntry{K03.11.rearrange_conditional}{Riemann rearrangement theorem}
\begin{theorem}
For a conditionally convergent real series, the positive terms sum to $+\infty$ and the negative terms sum to $-\infty$. A permutation can have any prescribed finite sum, can tend to either infinity, or can have any prescribed extended lower and upper limits $\alpha\leq\beta$.
\end{theorem}
\begin{proof}
Write $a_n=p_n-q_n$, $p_n=\max(a_n,0)$, $q_n=\max(-a_n,0)$. If one of $\sum p_n,\sum q_n$ were finite, convergence of $\sum a_n$ would make the other finite, contradicting $\sum\abs{a_n}=\infty$. Enumerate positive and negative terms in their original orders; both enumerations tend to zero.

For a finite target $s$, take unused positive terms until the partial sum first exceeds $s$, then unused negative terms until it first falls below $s$, and repeat. Each crossing exists by the infinite one-sided sums, and each overshoot is at most the last term used, tending to zero. All intermediate sums lie between the neighboring crossing sums up to these vanishing overshoots; hence they tend to $s$. Both lists are exhausted because infinitely many blocks of each sign are used. Insert zero terms successively, without affecting any sum.

For $+\infty$, at stage $j$ use the next negative term and then at least the next positive term, continuing positively until the sum exceeds $j$. The negative terms tend to zero, so even the minima of successive stages tend to $+\infty$; each list is exhausted. Reverse signs for $-\infty$.

For $\alpha<\beta$, choose finite $u_j<v_j$ with $u_j\to\alpha$, $v_j\to\beta$ in the extended sense. Alternately take at least one unused positive term and continue until above $v_j$, then at least one unused negative term and continue until below $u_j$. The crossing errors tend to zero; all intermediate values lie between neighboring endpoints, so the lower and upper limits are $\alpha,\beta$. The same proof covers infinite endpoints. When $\alpha=\beta$ use the finite or infinite constructions already given. This constructs the permutation, not merely a grouping.
\end{proof}

\KEntry{K03.11.cauchy_product}{Cauchy products and Mertens' theorem}
\begin{theorem}
Put $c_n=\sum_{k=0}^n a_kb_{n-k}$. If $\sum a_n=A$ absolutely and $\sum b_n=B$ converges, then $\sum c_n=AB$. If both factors are absolutely convergent, so is their Cauchy product.
\end{theorem}
\begin{proof}
Writing $B_j=\sum_{k=0}^j b_k$ gives $C_N=\sum_{n=0}^Nc_n=\sum_{k=0}^N a_kB_{N-k}$. Thus
\[C_N-AB=\sum_{k=0}^Na_k(B_{N-k}-B)-B\sum_{k>N}a_k.\]
The second term tends to zero. For the first, $\abs{B_j-B}\leq M$; choose $m$ with $M\sum_{k>m}\abs{a_k}<\eps$, and then make each of the finitely many factors $B_{N-k}-B$, $k\leq m$, small. This proves convergence. In the absolute case,
\[\sum_{n=0}^N\abs{c_n}\leq\sum_{k+j\leq N}\abs{a_k}\abs{b_j}\leq\Big(\sum_k\abs{a_k}\Big)\Big(\sum_j\abs{b_j}\Big),\]
so the product is absolutely convergent.
\end{proof}

\KEntry{K03.11.product_failure}{Two conditional factors can have a divergent product}
\begin{example}
Let $a_n=b_n=(-1)^n/\sqrt{n+1}$, $n\geq0$. Each series converges by the alternating test. Its Cauchy product has
\[\abs{c_n}=\sum_{k=0}^n\frac1{\sqrt{(k+1)(n-k+1)}}\geq\frac{2(n+1)}{n+2}\geq1,\]
because $(k+1)(n-k+1)\leq(n+2)^2/4$. Thus $c_n\not\to0$ and the product diverges. The factors are conditional, since $1/\sqrt{n+1}\geq1/(n+1)$.
\end{example}

\KEntry{K03.11.double_series}{Double sums and summation order}
\begin{theorem}
For nonnegative $a_{ij}$, all iterated sums and every enumeration of $\N_0^2$ have the same extended sum, the supremum over finite subsets. If $\sum_{i,j}\abs{a_{ij}}<\infty$, all such scalar sums exist and agree.
\end{theorem}
\begin{proof}
The nonnegative proof is exactly the finite-subset argument in \KRef{K03.11.nonnegative}, with rows or columns as blocks. In the absolute case choose a finite rectangle outside which the sum of magnitudes is $<\eps$ (a finite subset achieving this is contained in a rectangle). Any two partial sums containing that rectangle differ by $<2\eps$, proving existence and agreement; each iterated tail obeys the same bound.
For failure without the hypothesis, set $a_{ij}=1$ if $j=i$, $a_{ij}=-1$ if $j=i+1$, and zero otherwise. Every row sums to zero, whereas the first column sums to $1$ and every later column to zero. Thus the two iterated sums are $0$ and $1$.
\end{proof}

\KEntry{K03.11.products}{Infinite products and logarithms}
\begin{definition}
For nonzero scalar factors $v_n$, convergence of $\prod_{n\geq0}v_n$ means that $P_N=\prod_{n=0}^Nv_n$ tends to a nonzero scalar. A zero factor, or a nonzero-factor product whose limit is zero, is recorded separately, not called a convergent nonzero product. A convergent product necessarily has $v_n=P_n/P_{n-1}\to1$.
\end{definition}
\begin{theorem}
For positive factors, convergence to a positive nonzero product is equivalent to convergence of $\sum\log v_n$.
\end{theorem}
\begin{proof}
This is a downstream logarithmic characterization using \KRef{K03.07.log_construction} and \KRef{K03.07.exp_addition}: $\log P_N=\sum_{n=0}^N\log v_n$. Continuity of log proves one implication and continuity and positivity of exp prove the other. Finite nonzero prefixes simply multiply the limit. The example $v_n=1/2$ has limit zero and shows why the nonzero convention matters.
\end{proof}

\KEntry{K03.11.product_absolute}{Absolutely convergent products without logarithms}
\begin{theorem}
If $u_n\in\C$, $\sum\abs{u_n}<\infty$, and $1+u_n\ne0$ for all $n$, then $\prod(1+u_n)$ converges to a nonzero value.
\end{theorem}
\begin{proof}
For finitely many $t_j\geq0$ with $\sum t_j=T<1$, expansion gives $\prod(1+t_j)\leq\sum_{k\geq0}T^k=(1-T)^{-1}$: each elementary symmetric sum is at most $T^k$. Also
\[\abs{\prod(1+u_j)-1}\leq\prod(1+\abs{u_j})-1\leq T/(1-T).\]
Choose a tail with total magnitude $<1/2$. Its partial products, and hence the full partial products, are uniformly bounded; the displayed inequality for later tails makes $P_N$ Cauchy. Finally the elementary inequality $\prod(1-t_j)\geq1-\sum t_j$ (proved inductively) gives the modulus of every such tail product at least $1-T>1/2$. Multiplication by the fixed nonzero prefix bounds $\abs{P_N}$ away from zero. Its limit is therefore nonzero.
\end{proof}

\KEntry{K03.11.cesaro}{Cesaro summation}
\begin{definition}
With $s_N=\sum_{n=0}^Na_n$, the Cesaro means are $\sigma_N=(N+1)^{-1}\sum_{j=0}^Ns_j$. A finite limit of these means is the Cesaro sum.
\end{definition}
\begin{proof}
If $s_j\to S$, split the average of $s_j-S$ at a fixed $m$: the finite prefix divided by $N+1$ tends to zero, and the remaining average is bounded by $\sup_{j>m}\abs{s_j-S}$. Thus $\sigma_N\to S$. For $a_n=(-1)^n$, $s_N$ alternates between $1$ and $0$, so $\sigma_N\to1/2$ although the series diverges. Summability and ordinary convergence are distinct.
\end{proof}

\KEntry{K03.11.acceleration}{Remainders and a proved telescoping correction}
\begin{proposition}
An alternating remainder is bounded by its first omitted term. For a nonnegative decreasing $f$ on $[1,\infty)$ with convergent improper integral,
\[\int_{N+1}^\infty f(x)\dd x\leq\sum_{n>N}f(n)\leq\int_N^\infty f(x)\dd x.\]
For the reciprocal-square sum, a concrete correction is
\[\sum_{n=1}^\infty\frac1{n^2}=2+\sum_{n=1}^\infty\left(\frac1{n^2}-\frac1{n-1/2}+\frac1{n+1/2}\right),\]
and after $N$ corrected terms the remaining error has magnitude at most $1/(9N^3)$.
\end{proposition}
\begin{proof}
The first assertion is \KRef{K03.11.alternating}. The integral inequalities follow from $f(n+1)\leq f(x)\leq f(n)$ on $[n,n+1]$, summed on finite truncations before passing to the improper limit (\KRef{K03.10.integral_test}). For the correction,
\[\frac1{n-1/2}-\frac1{n+1/2}=\frac1{n^2-1/4},\qquad
\frac1{n^2}-\frac1{n^2-1/4}=-\frac1{4n^2(n^2-1/4)}.\]
The telescoping sum of the left difference through $N$ is $2-1/(N+1/2)\to2$. Since $n^2-1/4\geq3n^2/4$, the correction magnitude is at most $1/(3n^4)$. Its tail is at most $\frac13\int_N^\infty x^{-4}\dd x=1/(9N^3)$, proving both the identity and its useful error bound.
\end{proof}
